\documentclass[conference]{IEEEtran}

\usepackage{amsmath}
\usepackage{booktabs}
\usepackage{graphicx}
\usepackage{url}

\title{Hardware-Assisted Online Softmax Merge for Attention State Reuse on Spatz}

\author{
  \IEEEauthorblockN{Wenxuan Tang}
  \IEEEauthorblockA{Local Spatz Prototype Study}
}

\begin{document}

\maketitle

\begin{abstract}
This is the planned paper B skeleton. Unlike paper A, this manuscript is
reserved for a complete online softmax merge datapath and attention-like
evaluation. It should only claim complete mixed-scalar support after the RTL
datapath implements exponential scaling, weighted accumulation, and
normalization, and after the benchmark turns the current full-reference probe
from an expected-error case into an output comparison.
\end{abstract}

\section{Introduction}

Paper B extends the cluster-local merge-update platform from paper A toward a
complete online softmax merge engine for attention state reuse. The central
claim should be end-to-end or attention-like value, not only isolated
restricted-semantics microbenchmark speedup.

\section{Background}

This section should cover online softmax, attention tiling, state reuse, and
the Spatz/Snitch execution substrate~\cite{online-softmax2018,flashattention2022,spatz2023,snitch2020}.

\section{Architecture}

Describe the complete datapath:

\begin{align}
m_{new} &= \max(m_{old}, m_{tile}), \\
l_{new} &= l_{old}e^{m_{old}-m_{new}} +
           l_{tile}e^{m_{tile}-m_{new}}, \\
O_{new} &= O_{old}
          \frac{l_{old}e^{m_{old}-m_{new}}}{l_{new}} +
          O_{tile}
          \frac{l_{tile}e^{m_{tile}-m_{new}}}{l_{new}}.
\end{align}

This section should explicitly discuss `exp`, scaling, reciprocal/division,
pipeline latency, and TCDM scheduling.

\section{Numeric Design}

Reserve this section for approximation choices, LUT sweeps, reciprocal
implementation, and error metrics.

\section{Evaluation}

Required evidence before this paper is ready:

\begin{itemize}
  \item Full mixed-scalar correctness against CPU reference.
  \item CPU versus engine microbenchmark with the complete datapath.
  \item Attention-like or AttnRes workload results.
  \item Error statistics and resource/cost data.
\end{itemize}

\section{Limitations}

Keep this section honest about unsupported precision modes, workload coverage,
and resource data quality.

\section{Conclusion}

Summarize the complete datapath result only after the required evidence exists.

\bibliographystyle{IEEEtran}
\bibliography{common/bib/refs}

\end{document}
